% Abstract
\vspace{1in}
\chapter*{Abstract}

IT procurement covers the purchase, licensing, renewal and risk assessment of hardware, software and SaaS services. In many organisations it is still handled through e-mail, spreadsheets and separate tools, which leads to slow approvals, duplicate payments, unused licences, missed contract renewals and weak visibility of vendor risk. Existing techniques such as OCR, rule-based extraction, approval workflows and machine learning usually address one stage of the procurement cycle in isolation.

This report presents an Agentic AI-based IT Procurement Intelligence Platform built as an event-driven system of four cooperating agents: Document and Vendor Intelligence, Approval and Inventory Intelligence, Contract and Risk Analysis, and a Platform agent that provides authentication, business rules, notifications and a conversational assistant. Each agent is an independent FastAPI service that reacts to events on a Kafka-compatible bus (Redpanda), decides with rules or trained models, and publishes its result through a transactional outbox; long-running steps run as Temporal workflows, and decisions that are uncertain or sensitive are handed to a person.

Uploads are scanned with ClamAV and read with pdfplumber, Docling or PaddleOCR. A TF-IDF and logistic-regression classifier, regular-expression field extraction and an optional LayoutLMv3 cross-check produce a confidence-scored record, and records below a confidence of 0.8 go to a review queue. Vendors are matched by normalised fuzzy name, onboarding validates GSTIN check digits and IFSC codes, bank details require a second person's verification, and duplicate invoices are flagged. Requests follow spend-tier approval chains with SLA escalation and Redis-locked stock reservation; an Isolation Forest with SHAP explanations proposes licence reclaims after a grace period. Contracts are generated from templates, signed in the application, and read by a Legal-BERT model fine-tuned on CUAD, and vendors carry a Random Forest risk score with hard overrides for a cancelled GSTIN or a sanctions match, monitored for drift with the Population Stability Index.

On synthetic data, the shipped classifier reached a macro-F1 of 1.0 against 0.912 for a keyword baseline; Legal-BERT reached F1 scores of 0.85, 0.80 and 0.91 on three clause types of the CUAD test split, against 0.50, 0.19 and 0.02 for rules; the licence-anomaly model reached a ROC-AUC of 0.954; and the risk model reached a ROC-AUC of 0.715 against a ceiling of 0.739 set by its label noise. Field extraction was complete on text PDFs but reached only 49\% on photographed images. The implementation passes 823 service unit tests and four end-to-end suites through the gateway. All datasets are synthetic, so these figures compare methods rather than predict performance on an organisation's real documents.

\vspace{0.5cm}
\noindent\textbf{Keywords:} agentic AI, IT procurement, event-driven architecture, document intelligence, human-in-the-loop, vendor risk, anomaly detection, contract analysis.
